\chapter{Writer y backpressure}
\label{ch:writer}

\begin{note}
Stub. El task writer es el único punto de mutación para el
\Graph{} durante ingesta streaming. Un tratamiento completo
(preempción por enriquecimiento, contabilidad de quota,
overflow a disco) se difiere.
\end{note}

\section{Forma del pipeline}

\begin{verbatim}
LogSource -> raw_channel -> ParserTask -> parsed_channel -> GraphWriter -> Graph
\end{verbatim}

Ambos canales son colas \code{tokio::sync::mpsc} acotadas
(capacidades default 4 y 8 respectivamente, configurables por
fuente). Un canal lleno hace backpressure sobre el stage
aguas arriba; una fuente cuyo poll no encuentra lugar en
\code{raw\_channel} se bloquea en \code{send().await} hasta que
el writer haya drenado lo suficiente como para avanzar.

\section{Módulos}

\begin{itemize}
    \item \filepath{graph\_hunter\_core/src/ingest/batch.rs} ---
        unidad de trabajo \code{ParsedBatch} (triples +
        source\_id + dataset\_id + prioridad + watermark).
    \item \filepath{graph\_hunter\_core/src/ingest/parser\_task.rs}
        --- \code{ParserTask} consume \code{raw\_channel}, llama
        al \code{LogParser} específico del formato y emite
        \code{ParsedBatch}.
    \item \filepath{graph\_hunter\_core/src/ingest/writer.rs} ---
        \code{GraphWriter}: task singleton que drena
        \code{parsed\_channel}, toma el write lock del grafo por
        un batch a la vez, y lo libera entre batches para que
        los lectores puedan progresar.
    \item \filepath{graph\_hunter\_core/src/ingest/source\_poller.rs}
        --- \code{SourcePoller}: loop con timer por fuente que
        maneja \code{LogSource::poll\_since}.
    \item \filepath{graph\_hunter\_core/src/ingest/enrichment.rs}
        --- batches de enriquecimiento on-demand, preemptados
        por encima de la ingesta bulk.
    \item \filepath{graph\_hunter\_core/src/ingest/channels.rs}
        --- constructores de canales acotados.
    \item \filepath{graph\_hunter\_core/src/ingest/overflow.rs}
        --- serializa \code{ParsedBatch}es pendientes a disco en
        cancelación; replay en reinicio (Capítulo~\ref{ch:overflow}).
\end{itemize}

\section{Contención de locks}

El writer toma \code{graph.write()} durante una sola llamada a
\code{insert\_triples\_chunk} y la libera entre batches. Los
lectores (hunts, queries analíticas, queries de enriquecimiento)
toman \code{graph.read()} y coexisten entre sí. Los tamaños de
batch típicos son 1--5~K triples, así que un hold de write dura
milisegundos --- los lectores retroceden una vez por batch, lo
cual es barato.

\begin{invariant}[Unicidad del writer]
Cada sesión tiene a lo sumo un task \code{GraphWriter} activo en
un momento dado. El \code{SessionStore} registra el
\code{JoinHandle} del writer cuando arranca el task y rechaza
arrancar un segundo; esto descarta totalmente races
writer-writer.
\end{invariant}

\begin{inthecode}
\filepath{graph\_hunter\_core/src/ingest/writer.rs} --- loop
principal del writer.
\end{inthecode}
